% Generated from docs/REPORT_DRAFT.md; edit the manuscript then rerun the converter.
\section{Lệnh tái lập}

\begin{Verbatim}[fontsize=\small,breaklines=true,frame=single]
make setup
make data DATA_COUNT=2000
make run
make test
make experiments
make semantic-reasoning
make sparql-check
make evaluation-corpora
make neo4j-benchmark
make relational-benchmark
make evidence-summary
\end{Verbatim}

\section{Nguyên tắc phản biện}

\begin{itemize}
\item Không nói KG luôn nhanh hơn SQL.
\item Không gọi Cypher rule là OWL reasoning.
\item Không gọi silver corpus là gold/human evaluation.
\item Không gọi item-to-item similarity là personalization hoặc xác suất yêu thích.
\item Không nói Qwen trả lời; Qwen chỉ lập QueryPlan.
\item Không gọi hệ thống open-domain QA.
\item Không suy rộng benchmark 2.000 phim thành scalability result.
\end{itemize}

\section{Ma trận truy vết tiêu chí báo cáo}

\begin{longtable}{|>{\raggedright\arraybackslash}p{0.22\textwidth}|>{\raggedright\arraybackslash}p{0.22\textwidth}|>{\raggedright\arraybackslash}p{0.22\textwidth}|>{\raggedright\arraybackslash}p{0.22\textwidth}|}
\caption{Tổng hợp ma trận truy vết tiêu chí báo cáo}\label{tab:app-1} \\
\hline
\textbf{STT} & \textbf{Tiêu chí} & \textbf{Vị trí nội dung} & \textbf{Loại minh chứng} \\ \hline
\endfirsthead\hline \textbf{STT} & \textbf{Tiêu chí} & \textbf{Vị trí nội dung} & \textbf{Loại minh chứng} \\ \hline\endhead
1 & Bài toán, phạm vi, mục tiêu và lý do chọn KG & Chương 1; Mục 2.5 & lập luận và research question \\ \hline
2 & Survey và công trình gần đây & Chương 3 & protocol, 8 công trình, synthesis table, research gap \\ \hline
3 & So sánh mô hình thay thế & Mục 2.5; Chương 9 & bảng SQL/document/RDF/property graph và baseline \\ \hline
4 & Biểu diễn tri thức/ngữ nghĩa & Mục 2.1–2.4; Chương 5 & ontology, triple/quad, namespace, OWA/CWA \\ \hline
5 & Entailment và luật suy diễn & Mục 2.3; Chương 7 & semantic entailment và rule materialization \\ \hline
6 & Chất lượng tri thức và entity resolution & Mục 2.6, 2.9; Chương 6, 9 & policy, coverage, structural/semantic validation \\ \hline
7 & Công cụ, phiên bản và triển khai & Chương 6, 8; Phụ lục A & rationale, cấu hình và lệnh tái lập \\ \hline
8 & Dataset, nguồn và tiền xử lý & Mục 6.2–6.10 & provenance, terms, sampling, manifest và checksum \\ \hline
9 & Ontology/schema và mapping & Chương 5 & data dictionary và ánh xạ RDF–Neo4j \\ \hline
10 & Cài đặt graph/reasoner & Mục 6.7; 7.2–7.6 & import order, materialization và validation \\ \hline
11 & Ngôn ngữ truy vấn đặc thù & Mục 7.1, 7.4 & catalog Cypher/SPARQL và semantic path \\ \hline
12 & CRUD và nghiệp vụ & Mục 4.4–4.5; 7.1 & competency question và thao tác tham số hóa \\ \hline
13 & Inference-enabled query & Mục 7.2–7.6 & fact ẩn và evidence path \\ \hline
14 & Tiêu chí đánh giá & Mục 2.8; Chương 9 & metric, protocol và validity \\ \hline
15 & Benchmark, bảng, biểu đồ và baseline & Mục 9.2–9.7 & số liệu, protocol và placeholder hình evidence \\ \hline
16 & Bàn luận, hạn chế và cải tiến & Mục 9.8–9.12; Chương 10 & error analysis và roadmap ưu tiên \\ \hline
17 & Ứng dụng nghiệp vụ & Chương 4, 8 & stakeholder, QA/recommendation, API/UI và evidence \\ \hline
18 & Chất lượng báo cáo & toàn bộ báo cáo & bố cục, hình/bảng, BibTeX và phụ lục truy vết \\ \hline
\end{longtable}

Các tiêu chí phụ thuộc chương trình chạy được phải được đối chiếu với artifact
cuối; ma trận này chỉ bảo đảm người chấm tìm được nội dung và loại bằng chứng,
không thay thế kết quả thực nghiệm.

\section{Ngân hàng câu hỏi phản biện và trả lời}

\subsection{D.1. Vì sao không dùng CSDL quan hệ?}

CSDL quan hệ vẫn lưu được dữ liệu và là baseline hợp lệ. Neo4j được chọn vì
workload chính là pattern nhiều bước, shortest path và dựng evidence từ quan hệ.
Nếu workload chuyển thành giao dịch theo khóa hoặc aggregate cố định, hệ quan hệ
có thể là lựa chọn đơn giản hơn. Báo cáo không tuyên bố graph luôn nhanh hơn SQL.

\subsection{D.2. Vì sao cần đồng thời Neo4j và RDF/OWL?}

Neo4j là operational store thuận tiện cho property trên edge và traversal của
ứng dụng. RDF/OWL là standards view cho IRI, interoperability và entailment.
Hai biểu diễn phục vụ hai mục tiêu; mapping ở Chương 5 chỉ rõ phần tương đương và
phần mất mát như edge property cần reification/RDF-star.

\subsection{D.3. Cypher rule có phải OWL reasoning không?}

Không. \texttt{CO\_STARRED\_WITH} là luật nghiệp vụ closed-world được materialize từ hai
edge \texttt{ACTED\_IN} và lưu evidence. Semantic materializer xử lý tập con RDFS/OWL-RL
như domain, range, inverse và symmetric property. Báo cáo tách hai nhánh để không
đánh đồng procedural rule với ontology entailment.

\subsection{D.4. Tại sao không dùng tên làm khóa thực thể?}

Tên có thể trùng, đổi hoặc khác cách viết. Stable source ID giữ identity và cho
phép constraint. Fuzzy name chỉ dùng để liên kết truy vấn hoặc fallback có log;
trường hợp mơ hồ bị từ chối thay vì merge tự động.

\subsection{D.5. Vì sao rating TMDB và IMDb không gộp lại?}

Hai nguồn có population, thời điểm cập nhật và số phiếu khác nhau. Gộp làm mất
provenance. Hệ thống giữ hai thuộc tính và công khai chính sách \texttt{coalesce} khi
một use case cần một giá trị ưu tiên.

\subsection{D.6. Fact suy ra được kiểm chứng như thế nào?}

Mỗi \texttt{CO\_STARRED\_WITH} giữ \texttt{movie\_count} và \texttt{evidence\_movie\_ids}. Validation kiểm
tra phải tồn tại ít nhất một phim mà cả hai Person đều có \texttt{ACTED\_IN}; nếu tiền đề
không còn, fact suy ra bị coi là không được hỗ trợ.

\subsection{D.7. LLM có thể hallucinate Cypher hoặc câu trả lời không?}

LLM không được phép viết Cypher hoặc trả lời từ kiến thức riêng. Nó chỉ sinh JSON
theo \texttt{QueryPlan} schema. Entity linker ánh xạ slot, compiler whitelist sinh Cypher
tham số hóa, Neo4j trả fact và evidence. Kế hoạch sai schema hoặc ngoài intent bị
từ chối/fallback.

\subsection{D.8. Recommendation có phải cá nhân hóa không?}

Không. Đây là item-to-item similarity vì hệ thống không lưu user history/profile.
Điểm là tổng đóng góp IDF-weighted của feature chung, không phải xác suất người
dùng thích phim. Explanation faithful vì được tạo từ chính feature đã cộng điểm.

\subsection{D.9. Vì sao metric silver cao chưa chứng minh hệ thống hoàn hảo?}

Silver label được sinh từ quy tắc hoặc graph và có thể chia sẻ lỗi với hệ thống.
Do đó P/R/F1 hoặc precision cao chủ yếu chứng minh consistency với corpus, chưa
thay cho đánh giá gold độc lập. Báo cáo yêu cầu reviewer/adjudication trước khi
nâng claim.

\subsection{D.10. Benchmark hiện tại chứng minh điều gì?}

Nó mô tả latency của workload, snapshot, máy và protocol đã nêu. Một warm-up và
100 iteration hỗ trợ thống kê median/p95, nhưng một scale 2.000 phim không chứng
minh scalability. So sánh SQLite chỉ hợp lệ khi cùng snapshot, máy và cache policy.

\subsection{D.11. Những bias dữ liệu quan trọng nhất là gì?}

Popular-page sampling thiên về phim phổ biến; cast limit làm mất vai phụ; IMDb
match chỉ tồn tại khi TMDB có external ID; rating phản ánh nhóm người đã bình
chọn. Vì vậy coverage và recommendation không được suy rộng sang toàn bộ điện
ảnh hoặc mọi nhóm người dùng.

\subsection{D.12. Đóng góp quan trọng nhất của đề tài là gì?}

Đóng góp không nằm ở một thuật toán state-of-the-art mà ở workflow đầu cuối có
claim discipline: identity/provenance rõ, hai standards view, suy diễn có evidence,
execution surface của LLM được kiểm soát và evaluation phân biệt silver/gold.

\subsection{D.13. Nếu có thêm thời gian, nên mở rộng gì trước?}

Ưu tiên human review/adjudication, QA corpus lớn hơn và benchmark nhiều snapshot
thật; sau đó mới mở rộng Wikidata/Award, temporal model, vector retrieval hoặc
embedding. Thứ tự này cải thiện validity trước khi tăng độ phức tạp mô hình.

\section{Checklist trước khi xuất bản PDF}

\begin{itemize}
\item Điền đủ thông tin hành chính trên trang bìa và kiểm tra tên học phần.
\item Chạy lại toàn bộ thực nghiệm trên cùng snapshot rồi đồng bộ bảng số liệu.
\item Thay placeholder bằng hình thật đúng tên file; kiểm tra caption và khả năng đọc.
\item Kiểm tra danh mục hình, bảng, mục lục và số trang sau hai lần biên dịch.
\item Xác nhận mọi citation xuất hiện trong bibliography và mọi nguồn đã dùng được trích dẫn.
\item Kiểm tra bảng dài không tràn lề, code không vượt trang và hyperlink hoạt động.
\item Ghi số trang minh chứng vào bản checklist sau khi PDF đã khóa trang.
\item Không xóa phần limitation hoặc đổi silver result thành gold claim.
\end{itemize}
